\section{Development Results (2005--2021)}
\label{sec:dev}

\begin{table*}[t]
\caption{Development results, 2005--2021 (4,275 signal dates, 17 yearly folds). The portfolio holds 20 stocks against the equal-weight tradable universe; returns are net of 5\,bp per side unless stated. The DSR accounts for all 70 registered trials. \emph{Selected}: one LightGBM setting chosen each year; \emph{grid average}: 12 settings averaged without selection.}
\label{tab:dev}
\footnotesize
\setlength{\tabcolsep}{5pt}
\begin{tabular}{@{}lrrrrrrr@{}}
\toprule
Model & \ric{} & $t$ & Net active (\%/yr) & IR & IR at 10\,bp & Turnover (\%/yr) & DSR \\
\midrule
\tablefile{tables/dev_results.tex}
\bottomrule
\end{tabular}
\end{table*}

\begin{table*}[t]
\caption{Paired differences on the development period, with block-bootstrap intervals (21-day blocks, 2,000 draws). Bold intervals exclude zero; they are not corrected for multiple comparisons.}
\label{tab:dev-gaps}
\footnotesize
\begin{tabular}{@{}lll@{}}
\toprule
Comparison (A minus B) & \ric{} difference [95\% CI] & Net return difference, \%/yr [95\% CI] \\
\midrule
\tablefile{tables/dev_gaps.tex}
\bottomrule
\end{tabular}
\end{table*}

\begin{figure}[t]
\centering
% Cumulative net active return, 2005-2025. Data: figures/data/cumulative_active.tsv (export_data.py).
\begin{tikzpicture}
\begin{axis}[
  width=0.93\linewidth, height=5.3cm,
  xmin=2005, xmax=2026, ymin=-60, ymax=180,
  xtick={2005,2009,2013,2017,2021,2025},
  ytick={-50,0,50,100,150},
  x tick label style={/pgf/number format/1000 sep={}},
  tick label style={font=\footnotesize, text=black!75},
  ylabel={Cumulative net active (\%)}, ylabel style={font=\footnotesize},
  ymajorgrids, grid style={black!9},
  axis line style={black!35}, tick style={black!35}, tick align=outside,
  axis x line*=bottom, axis y line*=left,
  legend style={font=\footnotesize, draw=none, fill=none, at={(0.5,-0.14)}, anchor=north, column sep=5pt},
  legend columns=3, legend cell align=left,
]
\fill[black!6] (axis cs:2022,-60) rectangle (axis cs:2026,180);
\node[font=\footnotesize, text=black!60, anchor=north] at (axis cs:2024,178) {final test};
\draw[black!45] (axis cs:2005,0) -- (axis cs:2026,0);
\addplot[cCand, line width=1.2pt] table[x=t, y=candidate, col sep=tab] {figures/data/cumulative_active.tsv};
\addlegendentry{Primary model}
\addplot[cRev, line width=0.9pt, dashed] table[x=t, y=reversal, col sep=tab] {figures/data/cumulative_active.tsv};
\addlegendentry{Reversal}
\addplot[cLgb, line width=0.9pt, dash dot] table[x=t, y=lightgbm, col sep=tab] {figures/data/cumulative_active.tsv};
\addlegendentry{LightGBM}
\addplot[cTfmFt, line width=1.2pt] table[x=t, y=tfmft, col sep=tab] {figures/data/cumulative_active.tsv};
\addlegendentry{TimesFM FT}
\addplot[cTfmZs, line width=0.9pt, densely dotted] table[x=t, y=tfmzs, col sep=tab] {figures/data/cumulative_active.tsv};
\addlegendentry{TimesFM zero-shot}
\end{axis}
\end{tikzpicture}

\caption{Cumulative net active return against the equal-weight tradable universe (sum of daily active returns, 5\,bp per side). The shaded area is the final test, where the portfolios restart from cash. TimesFM FT is the 10k-step fine-tuned model.}
\Description{Line chart of cumulative net active return from 2005 to 2025. The primary model and reversal end near 140 and 130 percent, LightGBM and fine-tuned TimesFM near 80 and 95 percent, and zero-shot TimesFM declines to about minus 40 percent. Reversal gains most of its lead in 2008 to 2010.}
\label{fig:cumulative}
\end{figure}

\begin{figure*}[t]
\centering
% Left: Rank IC by year. Right: net active return against cost. Data from export_data.py.
\pgfplotsset{
  paperaxis/.style={
    tick label style={font=\footnotesize, text=black!75},
    label style={font=\footnotesize},
    ymajorgrids, grid style={black!9},
    axis line style={black!35}, tick style={black!35}, tick align=outside,
    axis x line*=bottom, axis y line*=left,
    legend style={font=\footnotesize, draw=none, fill=white, fill opacity=0.85, text opacity=1},
    legend cell align=left,
  },
}
\begin{tikzpicture}
\begin{axis}[
  paperaxis, name=ic,
  width=0.47\textwidth, height=5.0cm,
  xmin=2004.5, xmax=2025.5, ymin=-0.045, ymax=0.10,
  xtick={2005,2009,2013,2017,2021,2025},
  x tick label style={/pgf/number format/1000 sep={}},
  yticklabel style={/pgf/number format/.cd, fixed, precision=2},
  ylabel={\ric{} (yearly mean)},
  legend style={at={(0.5,1.03)}, anchor=south, fill=none}, legend columns=3,
]
\fill[black!6] (axis cs:2021.5,-0.045) rectangle (axis cs:2025.5,0.10);
\node[font=\footnotesize, text=black!60, anchor=north] at (axis cs:2023.5,0.099) {final test};
\draw[black!45] (axis cs:2004.5,0) -- (axis cs:2025.5,0);
\addplot[cCand, line width=1.1pt, mark=*, mark size=1.5pt] table[x=year, y=candidate, col sep=tab] {figures/data/rank_ic_by_year.tsv};
\addlegendentry{Primary model}
\addplot[cTfmFt, line width=1.1pt, mark=triangle*, mark size=1.9pt] table[x=year, y=tfmft, col sep=tab] {figures/data/rank_ic_by_year.tsv};
\addlegendentry{TimesFM FT 10k}
\addplot[cTfmZs, line width=0.9pt, densely dotted, mark=square*, mark size=1.3pt, mark options={solid}] table[x=year, y=tfmzs, col sep=tab] {figures/data/rank_ic_by_year.tsv};
\addlegendentry{TimesFM zero-shot}
\end{axis}

\begin{axis}[
  paperaxis,
  at={($(ic.east)+(1.2cm,0)$)}, anchor=west,
  width=0.47\textwidth, height=5.0cm,
  xmin=1.5, xmax=10.5, ymin=-11, ymax=11,
  xtick={2,5,10},
  xlabel={Cost per side (bp)}, ylabel={Net active return (\%/yr)},
  ytick={-10,-5,0,5,10},
  legend style={at={(0.02,0.02)}, anchor=south west}, legend columns=2,
]
\draw[black!45] (axis cs:1.5,0) -- (axis cs:10.5,0);
\addplot[cCand, line width=1.1pt, mark=*, mark size=1.5pt] table[x=cost, y=candidate, col sep=tab] {figures/data/cost_sensitivity.tsv};
\addlegendentry{Primary (2,386\%)}
\addplot[cRev, line width=0.9pt, dashed, mark=diamond*, mark size=1.9pt, mark options={solid}] table[x=cost, y=reversal, col sep=tab] {figures/data/cost_sensitivity.tsv};
\addlegendentry{Reversal (4,942\%)}
\addplot[cLgb, line width=0.9pt, dash dot, mark=pentagon*, mark size=1.7pt, mark options={solid}] table[x=cost, y=lightgbm, col sep=tab] {figures/data/cost_sensitivity.tsv};
\addlegendentry{LightGBM (4,347\%)}
\addplot[cTfmFt, line width=1.1pt, mark=triangle*, mark size=1.9pt] table[x=cost, y=tfmft, col sep=tab] {figures/data/cost_sensitivity.tsv};
\addlegendentry{TimesFM FT (3,535\%)}
\addplot[cMom, line width=0.9pt, loosely dashed, mark=square*, mark size=1.3pt, mark options={solid}] table[x=cost, y=momentum, col sep=tab] {figures/data/cost_sensitivity.tsv};
\addlegendentry{Momentum (215\%)}
\end{axis}
\end{tikzpicture}

\caption{Left: yearly \ric{} of the primary model and TimesFM, with the final test shaded. Right: development net active return against the cost per side; the legend gives each model's annual turnover, which sets the slope of each line.}
\Description{Left panel: yearly Rank IC from 2005 to 2025 for the primary model, fine-tuned TimesFM and zero-shot TimesFM; the primary model is positive in most years and zero-shot TimesFM often negative. Right panel: net active return at costs of 2, 5 and 10 basis points per side; reversal falls from about 10 to 2 percent, LightGBM from 6.5 to minus 0.4 percent, the primary model from 8.2 to 4.4 percent, fine-tuned TimesFM from 6.4 to 0.8 percent, and momentum stays near 1 percent.}
\label{fig:ic-cost}
\end{figure*}

\paragraph{Baselines.}
\cref{tab:dev} reports the development results, \cref{tab:dev-gaps} the paired comparisons and \cref{fig:cumulative} the cumulative active returns.
All four simple signals have a positive \ric{} with $t>2$.
Reversal earns the largest net active return of any single model (+6.9\%/yr, IR 0.54), but it trades 49 times its capital a year and earned most of it in 2008--2010 (+26.6\%/yr, against +1.1\%/yr over 2011--2021).
LightGBM on market features has a higher \ric{} (0.0263) but a lower IR (0.33): its turnover costs 4.3\% a year at 5\,bp, and its IR turns negative at 10\,bp.

\paragraph{What improves the tabular model.}
Averaging over the hyperparameter grid significantly raises the \ric{} once accounting and analyst features are present (+0.0051), for each of three seeds (+0.0039 to +0.0052), but not with market features only (+0.0006).
With more features, choosing one setting a year on a noisy validation metric seems to add more variance than it removes.
Accounting and analyst features add about 0.003 of \ric{} each and 0.0056 together, an interval that just includes zero, and cut turnover by a quarter.
The four simple signals lower turnover to 2,386\%/yr and lift the IR from 0.54 to 0.64, a gain within noise.
The resulting primary model has the best IR among models declared in advance, and it replaces reversal under our rule through its \ric{} (+0.0142, interval [0.0010, 0.0275]).
Its Newey--West $t$ is 2.8 and its six-factor alpha +4.2\%/yr ($t=1.7$, market beta 1.24), but its DSR of 0.68 stays below 0.95.

\paragraph{TimesFM~3.}
Zero-shot forecasts show no ranking skill (\ric{} 0.0072 with returns, 0.0038 with log prices), their portfolios trail the benchmark, and as LightGBM features they add nothing measurable.
Fine-tuning changes the picture: after 1,000 steps the \ric{} rises to 0.0247, significantly above zero-shot for both the \ric{} (+0.0175) and the net return (+5.6\%/yr), but a turnover of 4,633\% keeps the IR at 0.27.
Training ten times longer leaves the \ric{} unchanged ($-0.0019$), although the training correlation of the last block rises from 0.007 to 0.037, a sign that 1,000 steps under-trained it.
The 10k-step seeds agree more (mean rank correlation 0.56 against 0.44), resemble reversal less (0.36 against 0.51) and trade less (3,535\%), so the IR rises to 0.39 without a significant gain.
Fine-tuned TimesFM stays below the primary model in \ric{} ($-0.0107$, interval excluding zero), and adding it to the primary model lowers the IR from 0.64 to 0.56.

\paragraph{A striking but fragile ensemble.}
Fine-tuned TimesFM combined with reversal gives the highest IR of all 70 trials (0.80), a six-factor alpha of +8.0\%/yr ($t=2.6$) and a DSR of 0.87.
Its gains over reversal (+2.5\%/yr) and over the primary model (+2.6\%/yr) are not significant, and it depends on crises: its active return reached +61.6\% in 2008, and its IR is 0.53 without 2008--2010, against 0.61 for the primary model.
It did not replace the primary model, but it was declared as a comparator for the final test.

\begin{table}[t]
\caption{Variants of the primary model on the development period, declared before their results were computed; none of them can replace it. Active return and turnover in \%/yr, at 5\,bp unless stated.}
\label{tab:robustness}
\footnotesize
\setlength{\tabcolsep}{3pt}
\begin{tabular}{@{}lrrrrr@{}}
\toprule
Variant & \ric{} & Active & IR & IR 10\,bp & Turnover \\
\midrule
\tablefile{tables/robustness.tex}
\bottomrule
\end{tabular}
\end{table}

\paragraph{Robustness and costs.}
Eleven variants of the primary model declared in advance (\cref{tab:robustness}) keep a \ric{} between 0.030 and 0.034.
Their IR ranges from 0.24 to 0.64 with a median of 0.50, so the primary model sits at the top of its neighbourhood.
Removing reversal costs the most (IR 0.24); an exit rank of 50 matches the IR with 40\% less turnover, but choosing it after the fact would be one more selection.
\cref{fig:ic-cost} shows that turnover, not \ric{}, decides how a model reacts to costs: at 10\,bp the primary model keeps +4.4\%/yr, against +2.0\% for reversal and $-0.4$\% for LightGBM.
Charging each stock's closing half-spread plus 1\,bp gives the primary model an IR of 0.72, so 5\,bp is conservative for these liquid stocks, before market impact.
